\section{Evaluation Plan}
\label{sec:experiments}

No neural or benchmark empirical result is asserted in this draft. The audited
Phase-1 NumPy suite and the controlled linear-mixture confirmatory checks in
Section~\ref{sec:controlled-results} are conditional regression and
falsification artifacts, not performance results; the remaining experiments
are designed to falsify or support the theory rather than to fill a benchmark
table after the fact.

\subsection{Research Questions and Predictions}

\begin{description}
  \item[RQ1: Does capacity predict forgetting?] Test whether the
  curvature-whitened spectral tail in Eq.~\eqref{eq:spectral-capacity} and the
  exact and scale-invariant approximate safe dimensions
  $\operatorname{dim}(\cN_t^{\mathrm{amb}})$ and
  $\operatorname{dim}(\cN_t^{\mathrm{amb},(\tau)})$ predict the onset and slope of
  forgetting as $R$ varies.
  \item[RQ2: Does allocation matter at fixed capacity?] At a fixed deployed rank
  $R$, compare \method{} against implicit allocators that hold the same rank and
  buffer budget---dense sequential updating, reservoir/GSS/MIR replay, a
  gradient-projection defense, and the uniform-weight reduced-rank solve---and
  ask whether explicit reweighting changes worst-group retention at matched
  frequency-weighted mean.
  \item[RQ3: Is sequential accumulation conditional?] Measure the cross term in
  Eq.~\eqref{eq:forgetting-increment}. Streams with positive, near-zero, and
  negative cross terms should respectively show accumulation, diffusive growth,
  and cancellation or backward transfer.
  \item[RQ4: Does the local cost predict real loss?] Correlate the predicted
  historical change
  $\bq_{t-1}^\top\bh+\frac12\bh^\top\bA_{t-1}\bh$ with the observed replay and
  held-out historical loss change.
  \item[RQ5: Does the reweighting, not the buffer, drive any gain?] Ablate the
  minimax loop by fixing uniform weights (the uniform-weight reduced-rank solve),
  holding the buffer and closed-form solver identical, and test whether the
  worst-group gain survives.
\end{description}

\subsection{Controlled Temporal Mixture Streams}

The first controlled stream uses a single three-way financial-sentiment
instruction across Financial PhraseBank, FiQA-SA, and Twitter Financial News
Sentiment. Every window contains examples from multiple sources; only source
proportions and optional subdomains change. Because these datasets may share a
predictor, we treat this stream as a low-conflict covariate-shift control, not
as evidence that finite rank must become saturated. We instantiate four
schedules:
\begin{enumerate}
  \item \textbf{stationary mixture}, a negative control in which all
  $\pi_{t,c}$ are fixed;
  \item \textbf{smooth drift}, with slowly varying mixture weights;
  \item \textbf{abrupt recurrent drift}, in which earlier mixtures return; and
  \item \textbf{expanding support}, in which new sources or subdomains enter.
\end{enumerate}

The primary real-data capacity stress test is a mixed-skill instruction stream.
Each window mixes examples from several source families---for example
extraction, question answering, summarization, reasoning, and code---under one
causal-LM or supervised instruction loss. Mixture weights drift and support
expands, while source metadata is retained only for analysis. Before comparing
methods, a held-out pilot estimates the curvature-whitened rank of window
optimum or gradient innovations. We pre-register the diagnostic and report both
low-innovation and high-innovation streams; datasets are not selected after
viewing \method{} performance.

A synthetic linearized stream independently controls the rank, energy, and
mutual alignment of new curvature modes. It is the cleanest test of
Theorems~\ref{thm:spectral-capacity} and~\ref{thm:safe-subspace}, because the
ground-truth spectral tail and nullspace dimension are known.

\subsection{Chronological Language-Model Streams}

TemporalWiki and StreamingQA provide timestamped knowledge updates
\citep{jang2022temporalwiki,liska2022streamingqa}. We use one causal-LM or
instruction-tuning objective throughout and do not expose timestamps as task
IDs. Evaluation partitions questions into stable facts, new facts, updated
facts, and obsolete answers. A successful method should retain stable facts,
acquire new facts, answer updated facts with their current values, and avoid
persisting obsolete answers. TRACE or Super-NI sequences are included only as
supplemental task-incremental comparisons.

Chronological claims require a base model with a documented pretraining cutoff
before the evaluated stream, or a newly constructed post-cutoff slice. We log
the unadapted model's closed-book accuracy and answer likelihood as
contamination probes and report suspected memorized items separately. Fact
categories are generated from timestamped revision or answer changes and a
stratified sample is manually audited. Every updated fact receives two labels:
its value at the historical checkpoint and its current value. Frozen-truth
retention and current-truth adaptation are reported separately, so a correct
revision is not counted as forgetting.

\subsection{Baselines and Resource Matching}

\begin{table}[t]
\caption{Planned comparisons along separate control axes. No single baseline
can simultaneously match dense rank-$R$ activation and active rank $k<R$.
Oracles are reported separately and never described as online peers.}
\label{tab:baseline-plan}
\centering
\small
\begin{tabularx}{\linewidth}{
  >{\raggedright\arraybackslash}p{0.18\linewidth}
  X
  >{\raggedright\arraybackslash}p{0.18\linewidth}}
\toprule
Group & Method & Controlled quantity \\
\midrule
Capacity-matched & Sequential LoRA; dense rank-$R$ activation & Deployed rank
$R$, history bytes, tokens \\
Rank-and-buffer-matched & \method{}; uniform-weight reduced-rank solve;
reservoir/GSS/MIR replay; gradient projection & Deployed rank $R$, buffer
capacity $B$, and history bytes \\
History-memory-matched & Replay-LoRA; gradient projection; \method{} &
Total replay and curvature-sketch bytes \\
Mechanism & O-LoRA; SLICE; E$^2$-LoRA; ProCL; NSR; LiteLoRA & Full ledger,
including banks, routers, retrieval, and extra passes \\
Consolidation & Euclidean SVD; DELLA; TIES; curvature-weighted SVD & Secondary \\
Oracles & Offline joint training; growing rank; one adapter per window with
oracle selection; full fine-tuning & Upper/resource bounds \\
\bottomrule
\end{tabularx}
\end{table}

For every method we report total and active trainable parameters, deployed
adapter bytes, replay or sketch bytes, optimizer-state bytes, wall-clock time,
training FLOPs, and inference cost. Recent allocation methods can use additional
retrieval banks or input-conditioned gates, so parameter count alone is not a
sufficient fairness criterion. Capacity claims use the capacity-matched group;
allocation-efficiency claims use the rank-and-buffer-matched group; historical
memory claims use the history-memory-matched group. Full historical evaluation
sets are read only for reporting. The teacher estimate, the minimax reweighting,
and the kept-iterate selection may inspect only the bounded buffer already
counted in the ledger.

\subsection{Metrics, Statistical Protocol, and Diagnostics}

Primary performance metrics are average online performance, final average
performance, backward transfer, loss-based forgetting, worst-window retention,
stable-fact retention, updated-fact accuracy, obsolete-answer persistence, and
general-capability drift. Capacity diagnostics include effective rank,
curvature-whitened spectral tail,
$\operatorname{dim}(\cN_t^{\mathrm{amb}})$,
$\operatorname{dim}(\cN_t^{\mathrm{amb},(\tau)})$, the numerical rank of the
deployed operator, the minimax weight vector and buffered worst-window loss per
round, the kept-iterate index, and the distance of each arm's mean--tail point to
the offline frontier.

We run at least three seeds and target five for principal results. We report
means, standard deviations, paired bootstrap confidence intervals over matched
stream realizations, and per-stream win rates. ``Significant'' is reserved for
a pre-specified test with an uncertainty interval. Rank sweeps use
$R\in\{8,16,32,64\}$ and $k\in\{2,4,8\}$ where model size permits, together
with replay-memory sweeps. Exported checkpoints are audited to verify their
actual matrix rank; requested rank alone is not accepted as evidence.

\paragraph{Failure criteria.}
The capacity thesis is weakened if spectral tails and safe-space counts do not
predict forgetting better than update norm or time, if fixed-budget allocation
does not beat dense/random policies under matched resources, or if curvature
cost is uncorrelated with observed historical loss. The method claim fails if
its gains disappear after accounting for replay, optimizer state, or extra
training compute.

\subsection{Controlled-Mixture Confirmatory Results}
\label{sec:controlled-results}

The core theory-consistency checks (D1) have been reproduced both locally and on
the \texttt{alibaba-10} Linux host with matching gates; the rare-task bridge
(D2-A) and the nonlinear-transfer probe (D2-B) were run locally with
source-hash-recorded manifests. All are conditional confirmatory artifacts, not
neural or benchmark results. The \method{} evaluation (D2-C) is a controlled
linear-factored-proxy performance result under an internally pre-specified,
hash-recorded held-out kill test, and is explicitly not a real-LoRA,
language-model, or benchmark claim.

\paragraph{Expressivity predicts forgetting (D1).}
On a planted linear-LoRA stream with two disjoint window subspaces
($d=8$, $T=6$, rank $R=2$), the curvature-weighted Grassmann distance
$d_t$ (Proposition~\ref{prop:expressivity-metric}) is zero on an identical
occupied subspace and $d_t=\sqrt{2}$ between same-rank adapters aimed at
disjoint windows, confirming that equal rank does not imply expressive
equivalence. Early-window adapters retain early windows better and vice versa
($50.2$ vs.\ $295.2$ forgetting on late windows), so the distance predicts the
direction of retention. Across the three occupied-rank points probed here, the
effective capacity $\kappa_t$ (Proposition~\ref{prop:effective-capacity}) is
rank-correlated with forgetting ($+1$); we note that the effective-rank and
update-norm baselines are equally perfectly rank-correlated ($-1$) on these
three points, so this small sweep establishes consistency with the theory, not
separation from the baselines, which would require a denser sweep. As an empirical comparator, a full fine-tuning reference at its effective rank
$r_{\mathrm{eff}}=6$ lands at the same $\kappa_t$-forgetting point as the rank-6
LoRA adapter (forgetting difference $<10^{-14}$); we read this as a diagnostic
that full fine-tuning behaves like its effective, not nominal, rank on this
stream, not as a theorem unifying the two training regimes.

\paragraph{Expressivity predicts rare-task retention (D2-A).}
On the linear mixture of \citet{huang2026larger} ($K=8$ orthogonal tasks,
spectra $j^{-2}$, frequencies $k^{-1.5}$), the curvature-whitened $N$-th value
$\mu_N^{\mathcal{F}}$ of the common block reproduces their predicted critical
width: at $N^{\mathrm{crit}}=12$, $\mu_{12}^{\mathcal{F}}=0.0204\le\pi_r\lambda_r=0.0229$.
The tail sum $\kappa_t$ decreases with width, the distance from the occupied
common subspace to the rare direction decreases, and $\kappa_t$ correlates with
inverse retention (Spearman $0.82$). Crucially, a width-$N$ encoder optimized by
Adam reproduces the analytic onset: at $N=10<N^{\mathrm{crit}}$ the rare signal
is $2.5\times10^{-7}$, while at $N=14>N^{\mathrm{crit}}$ it reaches $0.73$. The
controlled, non-neural optimizer therefore respects the curvature-weighted
prediction.

\paragraph{The mechanism does not yet transfer to a minimal nonlinear model (D2-B, negative).}
We ran the same rare-task probe on a minimal nonlinear model---a two-layer ReLU
MLP with a width-$N$ bottleneck ($d_{\mathrm{in}}{=}32$, $K{=}8$ tasks), not a
Transformer---trained by Adam, along two capacity axes: a frozen-encoder adapter
rank sweep and a full-fine-tuning encoder width sweep. The curvature-weighted
prediction \emph{fails} to transfer: rare-task loss does not decrease
monotonically with rank ($2/3$ gates fail) or with width ($3/4$ gates fail); at
high width the rare task is still not learned, and the $\kappa_t$--loss rank
correlation is only $0.29$ (rank sweep) and $0.5$ (width sweep) rather than the
$1.0$ seen in the linear reductions. We report this negative result rather than
suppress it: the curvature-tail account is presently validated only in the
controlled linear and Kronecker-quadratic settings, and identifying what breaks
in the nonlinear regime (optimizer path dependence, feature learning, or a
mis-specified curvature surrogate) is required future work before any real-LoRA
or LLM claim. An OLMo-scale confirmation is therefore not yet warranted.

\paragraph{Explicit reallocation redistributes the fixed budget across sources (D2-C).}
The scale-driven mechanism of \citet{huang2026larger} raises capacity to reduce
interference; we instead hold capacity fixed and ask whether \emph{how} the
budget is allocated changes worst-group retention. The arena is a single global
linear teacher $\bT^\star$ in the capacity-limited regime---$K$ sources with
correlated input directions, so the rank-$R$ tail is redistributable across
sources rather than zero-sum, and a fixed budget $R{=}4{<}K$---the only one of
our regimes with a genuine rank-$R$ floor. Every arm is oracle-free (a function
of $(\bX,\bY)$ only, verified by a label-blindness hash), sees the same clocked
windows, and holds the same bounded buffer of $B{=}16$ windows, matching the
replay baselines' budget. Retention is measured on independent held-out batches:
$\mathrm{mean}$ is frequency-weighted source retention and $\mathrm{tail}$ is
worst-group retention. The two are distinct functionals: what
Corollary~\ref{cor:capacity-lower-bound} bounds is the \emph{aggregate}
frequency-weighted capacity error at rank $R$, which no allocation rule can
lower; worst-group retention is instead the object that reweighted reallocation
\emph{redistributes}---it selects which source the fixed aggregate floor is
charged to, not whether the floor exists. A closed-form weighted reduced-rank
solve computed offline gives the reachable frontier, whose best tail is $0.839$.

To prevent tuning to the outcome we fixed a kill criterion in committed source
\emph{before} running the held-out set---the evaluation script states the criterion,
the worlds, and a ``run once, report whatever comes out'' rule in its header, and
the method file freezes the exact configuration as its default---and ran it once.
We release SHA-256 hashes of these artifacts to pin them; the hashes were recorded
at release time and pin the released code rather than independently timestamping
that it predated the numbers, which instead rests on the committed script. The coverage split ($B_u{=}4$,
$B_g{=}12$) and every hyperparameter were selected on a disjoint development
family of teacher and stream seeds; the test set is twenty held-out worlds---five
new teacher geometries $\times$ four new stream seeds, a crossed $5\times4$
design---disjoint from the development family and from an earlier burned test set,
and never seen during design. The method is kept only if all three hold, paired
over the twenty worlds: (a)~\textbf{non-inferior mean} within a $0.01$ margin of
the strongest-mean baseline (excess-CVaR here); (b)~\textbf{strictly higher tail}
than the strongest risk-aware replay baseline (excess-CVaR), with a paired
confidence interval excluding zero; and (c)~\textbf{mechanism contribution},
tail strictly above the uniform-weight solve on the same coverage buffer---the
ablation that shares the buffer and closed-form solver but drops the
reweighting---under the same interval test. All decisions use the cluster-robust
two-way crossed $95\%$ interval for the teacher$\times$stream design (conservative
$\mathrm{df}=\min(n_t,n_s)-1=3$); the IID interval is reported only for
transparency and is never the decision object. If (c) fails, the minimax loop
adds nothing over ``buffer plus closed-form solve'' and the method line stops
regardless of (a) and (b).

The verdict is KEEP: all three criteria pass on the held-out set
(Table~\ref{tab:d2c}). \method{} reaches worst-group retention
$0.792\pm0.062$, the best of any arm and $94\%$ of the offline frontier's
$0.839$, at frequency-weighted mean $0.970\pm0.005$. The pure implicit
allocators---dense sequential updating, reservoir/GSS/MIR replay, and an A-GEM
gradient-projection defense---cluster at tail $\approx0.58$--$0.60$, the
implicit-allocation floor, while matching \method{} on mean to within $0.005$.
The three paired tests, all on the cluster-robust crossed interval:
(a)~mean vs.\ the strongest-mean baseline (excess-CVaR) is $+0.0027$ (cluster
$95\%$ CI $[-0.0026,+0.0079]$), inside the $-0.01$ margin; (b)~tail vs.\
excess-CVaR is $+0.1782$ (cluster $95\%$ CI $[+0.0696,+0.2869]$), excluding
zero; and (c)~tail vs.\ the uniform-weight solve on the same coverage buffer is
$+0.0945$ (cluster $95\%$ CI $[+0.0477,+0.1413]$), excluding zero.

The tail gain decomposes into two roughly additive mechanisms, each about
$+0.09$ over the $\approx0.60$ implicit floor. Reweighting alone
(\texttt{psr\_no\_coverage}: minimax reweighting on a single uniform reservoir at
the same budget, no coverage) reaches $0.701$, and coverage alone
(\texttt{buffer\_rrr\_uniform}: the coverage buffer with uniform weights) reaches
$0.698$; only their combination reaches $0.792$. We are explicit about which increments are
interval-strict. Criteria (b) and (c) are: reweighting given coverage clears the
interval (c), and the full method clears it against the replay bar (b). The
coverage increment on its own is \emph{not}: measured against
\texttt{psr\_no\_coverage} it is $+0.0918$ but its cluster interval
$[-0.0085,+0.1922]$ straddles zero (it wins $18$ of $20$ worlds), so we report
coverage as a jointly-necessary component of the combination (only coverage and
reweighting together reach $0.792$), not as an independently significant effect,
and we do not claim a formal super-additivity of the two increments. We are equally explicit that (b) is a paired-average
superiority, not uniform dominance: \method{} beats excess-CVaR on $19$ of $20$
worlds; the mechanism ablation (c) is cleaner still, with \method{} above the
uniform solve on all $20$.

\begin{table}[t]
\centering
\small
\caption{Hash-recorded held-out kill test for \method{} on the capacity-limited
regime, over twenty held-out worlds (five new teacher geometries $\times$ four new
stream seeds, a crossed $5\times4$ design; $R{=}4$, buffer $B{=}16$, oracle-free,
independent held-out retention). Tail is worst-group retention; ``tail/off'' is
the fraction of the offline weighted-RRR frontier ($0.839$). Arms are sorted by
tail. The pure implicit allocators (dense sequential, replay variants, A-GEM,
EWC) sit at the implicit-allocation floor ($\mathrm{tail}\approx0.58$--$0.60$).
The two ablations isolate the mechanisms: reweighting on a single reservoir
(reweight only) and uniform weights on the coverage buffer (coverage only) each
reach $\approx0.70$, and only their combination (\method{}) reaches $0.792$, at
matched mean. Criteria (a) vs.\ excess-CVaR, (b) vs.\ excess-CVaR, (c) vs.\ the
coverage-only ablation---all pass on the cluster-robust crossed interval (see
text).}
\label{tab:d2c}
\begin{tabular}{lccc}
\toprule
Arm & Mean & Tail (worst-group) & tail/off \\
\midrule
\textbf{\method{} (coverage $+$ reweight)} & $\mathbf{0.970\pm0.005}$ & $\mathbf{0.792\pm0.062}$ & $\mathbf{0.94}$ \\
\;\;reweight only, single reservoir (abl.)  & $0.970\pm0.005$ & $0.701\pm0.096$ & $0.83$ \\
\;\;coverage only, uniform weights (abl.)   & $0.973\pm0.006$ & $0.698\pm0.092$ & $0.83$ \\
\midrule
grad-defensive (neg.\ control)  & $0.968\pm0.008$ & $0.630\pm0.148$ & $0.75$ \\
excess-CVaR replay              & $0.968\pm0.007$ & $0.614\pm0.119$ & $0.73$ \\
A-GEM (grad.\ projection)       & $0.967\pm0.007$ & $0.600\pm0.150$ & $0.71$ \\
EWC-LoRA                        & $0.965\pm0.009$ & $0.596\pm0.163$ & $0.71$ \\
sequential (dense)              & $0.966\pm0.008$ & $0.594\pm0.179$ & $0.71$ \\
MIR replay                      & $0.960\pm0.010$ & $0.589\pm0.137$ & $0.70$ \\
GSS replay                      & $0.966\pm0.007$ & $0.585\pm0.152$ & $0.70$ \\
uniform replay                  & $0.965\pm0.008$ & $0.577\pm0.152$ & $0.69$ \\
\midrule
offline oracle (reachable)      & --- & $0.839$ & $1.00$ \\
\bottomrule
\end{tabular}
\end{table}

\paragraph{Off-floor robustness: no harm and graceful degradation.}
The redistribution gain should appear only where a rank-$R$ floor exists, and the
reallocation must not misbehave elsewhere. We therefore ran two off-floor checks
at the same frozen coverage config (three teachers $\times$ two streams each). On
the \emph{compressible} regime (no rank-$R$ floor: the teacher fits within the
budget) all arms are near-perfect and \method{} does no harm
($\mathrm{mean}/\mathrm{tail}=0.992/0.988$, tying the coverage-only solve's
$0.992/0.988$ and excess-CVaR's $0.992/0.988$). On the \emph{conflicting} regime
(an impossibility control where sources demand contradictory targets, so the tail
is largely unrecoverable in principle) every arm degrades sharply; \method{}
degrades gracefully, matching the coverage-only solve on mean ($0.530$ vs.\
$0.531$) and on tail ($0.237$ vs.\ $0.235$), and it does not fall below the
uniform solve or manufacture a spurious advantage---its small tail edge over
excess-CVaR ($0.237$ vs.\ $0.189$) is the same order as on the capacity floor and
is not claimed as a rescue. The reallocation helps at the floor and is inert or
safe away from it.

\paragraph{Scope of the \method{} result.}
D2-C is a controlled linear-factored-proxy result. It authorizes no claim about
real LoRA training, language models, or benchmark performance, and the tail gain
is a paired-average effect on held-out geometries rather than a per-world
guarantee. Combined with the negative transfer in D2-B, the honest boundary is
that explicit reallocation \emph{redistributes} the fixed budget to raise
worst-group retention at matched frequency-weighted mean---it does not and cannot
lower the aggregate capacity floor---\emph{in the linear regime where that floor
is exactly characterized}; a nonlinear and real-LoRA confirmation is required
future work.
